Facility location decisions in dense urban environments depend not only on
population size but also on network accessibility, territorial structure,
and socioeconomic heterogeneity. This paper presents a modular geospatial
optimization framework that integrates pedestrian-network distances derived
from OpenStreetMap, graph-based spectral clustering, multiple correspondence
analysis~(MCA), and a clustered \mbox{p-median} formulation with coverage,
separation, and demand-penalty constraints.
The methodological contribution is a complete, calibrated pipeline in which
every parameter is justified through multi-objective sensitivity
analysis over an 81-point grid rather than fixed \textit{a priori}.
A proportional-demand allocation rule replaces the original fixed-per-cluster
budget, reducing inequity across territories of different size.

The framework is validated across three urban contexts in Mexico that differ
substantially in density, infrastructure supply, and commercial maturity:
a dense student-area district (Tecnológico de Monterrey, Monterrey),
a less-developed city~(Villahermosa, Tabasco), and an industrial
corridor~(Estado de México).
In all three settings the clustered model is compared on equal terms against
four baselines evaluated under pedestrian-network distance:
a global \mbox{p-median} optimum, a \mbox{k-means} partitioning baseline,
an Euclidean-solve baseline, and a no-MCA variant.
The proposed model retains \SIrange{93.2}{99.9}{\percent} of the global
optimum coverage while reducing solver runtime by up to \SI{67}{\percent}
in larger instances.
Euclidean-distance models underestimate mean assignment distance by
\SIrange{10}{20}{\percent} relative to network evaluation, confirming
that pedestrian routing is an essential component of any urban
accessibility assessment.
The framework is modular and transferable to clinics, pharmacies,
service offices, and other location systems where
network accessibility and territorial heterogeneity are primary concerns.
